% Problem record op_2fba72464e1b4de4. This TeX file is derived from
% database/problems_json/op_2fba72464e1b4de4.json by scripts/sync-tex.mjs; edit the JSON record
% and rerun the script rather than editing this file.
\section{Minimal block length of permutation-invariant qubit codes}
\label{sec:op_2fba72464e1b4de4}

\paragraph{Problem.}
Is the least number of physical qubits on which a two-dimensional
permutation-invariant code corrects arbitrary errors on $t$ qubits equal to
$3t^2+3t+1$ for every integer $t\geq1$?  For $n\geq1$ let
$\operatorname{Sym}^n(\mathbb C^2)\subseteq(\mathbb C^2)^{\otimes n}$ denote the
symmetric subspace, consisting of the vectors fixed by every permutation of the
$n$ tensor factors; a code contained in it is fully permutation-invariant, with
every code vector, not merely the code projector, invariant under relabelling
the qubits.  A subspace $\mathcal C$ with orthogonal projector $P$ has distance
$d(\mathcal C)\geq d$ when $PEP$ is a scalar multiple of $P$ for every Pauli
operator $E$ acting nontrivially on fewer than $d$ qubits, and it corrects
arbitrary errors on $t$ qubits exactly when $d(\mathcal C)\geq 2t+1$.  Define
\begin{equation}
  n_{\min}^{\mathrm{PI}}(t):=\min\bigl\{n\in\mathbb N:\
    \exists\,\mathcal C\subseteq\operatorname{Sym}^n(\mathbb C^2),\
    \dim\mathcal C=2,\ d(\mathcal C)\geq 2t+1\bigr\}.
  \label{eq:2fba-nmin}
\end{equation}
The minimum in Eq.~\eqref{eq:2fba-nmin} ranges over all complex
two-dimensional subspaces of the symmetric subspace and over both even and odd
$n$; no coefficient ansatz, reality condition, or transversal-gate requirement
is imposed.  The question is whether
\begin{equation}
  n_{\min}^{\mathrm{PI}}(t)=3t^2+3t+1
  \qquad\text{for every integer } t\geq1 .
  \label{eq:2fba-conjecture}
\end{equation}
A complete answer proves Eq.~\eqref{eq:2fba-conjecture} or exhibits an integer
$t\geq1$ for which it fails.

\subsection*{Status}

Unsolved

\subsection*{Source}

Bond, Min\'a\v{r}, Ozols, Safavi-Naini, and Visnevskyi state
Eq.~\eqref{eq:2fba-conjecture} explicitly as Conjecture~3.1, on the basis of the
numerical searches of their Section~3.1 and Figs.~1 and~2
\sourcecite{ref:2fba-bond-et-al}{BMO+26}.  The quantity in
Eq.~\eqref{eq:2fba-nmin} extends the single-error question studied by Pollatsek
and Ruskai \sourcecite{ref:2fba-pollatsek-ruskai}{PR04}.

\subsection*{Progress}

\begin{itemize}
  \item Pollatsek and Ruskai constructed two seven-qubit permutation-invariant
  codes that correct every single-qubit error (Section~6.2, Eq.~(62)), so
  $n_{\min}^{\mathrm{PI}}(1)\leq7$.  Their necessary and sufficient conditions
  (Theorem~1) and their exclusion of a five-qubit code (Section~6.1) concern real
  coefficients, odd lengths, and a restricted code form, not every
  two-dimensional subspace of the symmetric subspace
  \sourcecite{ref:2fba-pollatsek-ruskai}{PR04}.

  \item Ouyang's gnu codes with $g=n=2t+1$ and scaling $u\geq1$ correct arbitrary
  errors on $t$ qubits (Theorem~4).  At $u=1$ their length is $(2t+1)^2$, so
  $n_{\min}^{\mathrm{PI}}(t)\leq(2t+1)^2$ for every $t$, and permutation
  invariance is compatible with correcting any fixed number of errors
  \sourcecite{ref:2fba-ouyang}{Ouy14}.

  \item Aydin, Alekseyev, and Barg define codes $Q_{g,m,\delta,\epsilon}$ of length
  $2gm+\delta+1$ and prove in Theorem~5.3 that they correct arbitrary errors on
  $t$ qubits whenever $m\geq t$, $\delta\geq 2t$, and either $g\geq 2t$ with
  $\epsilon=-1$ or $g\geq 2t+1$ with $\epsilon=+1$.  The choice
  $Q_{2t,t,2t,-}$ gives the explicit bound
  \begin{equation}
    n_{\min}^{\mathrm{PI}}(t)\leq 4t^2+2t+1 ,
    \label{eq:2fba-aab-bound}
  \end{equation}
  which matches Eq.~\eqref{eq:2fba-conjecture} only at $t=1$
  \sourcecite{ref:2fba-aab}{AAB24}.

  \item Section~7 of the same paper extends the Pollatsek--Ruskai conditions to
  every $t$ for real, odd-length codes of a fixed parity-separated form
  (Proposition~7.1).  For $t=2$ and $n=19$ these conditions are nine quadratic
  equations, and the authors report a real solution computed with the
  polynomial-system solver msolve, whose output is an interval for each
  variable containing an exact solution, cross-checked in Mathematica.  This
  establishes a nineteen-qubit permutation-invariant code of distance five, so
  $n_{\min}^{\mathrm{PI}}(2)\leq19$.  The accompanying remark that no shorter
  $t=2$ code exists refers to that ansatz and is stated without proof
  \sourcecite{ref:2fba-aab}{AAB24}.

  \item Bond and collaborators minimize a Knill--Laflamme cost function, accepting a
  solution when the cost is below $10^{-18}$ and the gradient below $10^{-20}$,
  with up to $1000$ random starts per length.  With unrestricted complex
  coefficients the first solutions appear at $n=7,19,37$ for $t=1,2,3$; within
  the Pollatsek--Ruskai family they appear at $n=7,19,37,61,91$ for
  $t=1,\dots,5$ (Section~3.1, Figs.~1 and~2).  They also note the quantum
  Singleton bound $n\geq 4t+1$.  Local numerical searches neither certify codes
  nor exclude shorter ones, so this is evidence for
  Eq.~\eqref{eq:2fba-conjecture} rather than a proof
  \sourcecite{ref:2fba-bond-et-al}{BMO+26}.

  \item In a different regime, Aydin, Albert, and Barg construct
  permutation-invariant codes whose local dimension $q$ equals the block length
  $N$, with code dimension $K=o(2^N)$ and distance $d=o(N/\log N)$
  (Proposition~VII.7 and Theorem~VII.8).  The growing local dimension is
  essential, so these codes do not bound the qubit quantity in
  Eq.~\eqref{eq:2fba-nmin} \sourcecite{ref:2fba-aab-convex}{AAB26}.

  \item Kubischta and Teixeira identify subspaces of $\operatorname{Sym}^n(\mathbb C^2)$
  with intrinsic codes in the spin-$n/2$ representation of $\mathrm{SU}(2)$,
  with intrinsic depth equal to qubit distance (Theorem~B1, Lemma~B1, and
  Proposition~B2 of version~2).  Their resulting linear program has a unique
  feasible enumerator for $n=7$, $K=2$, and distance three
  (Propositions~B5--B6), and Appendix~B-D states that the same program is
  infeasible at every shorter length, so no permutation-invariant $((n,2,3))$
  code exists for $n<7$; the routine infeasibility checks are described in
  Appendix~B-B but not written out.  With the seven-qubit codes this gives
  $n_{\min}^{\mathrm{PI}}(1)=7$ without any reality or ansatz restriction.  The
  result appears in the 26 August 2026 revision of a preprint
  \sourcecite{ref:2fba-kubischta-teixeira}{KT26}.

  \item Teixeira identifies the intrinsic MacWilliams matrix of permutation-invariant
  qudit codes with a Racah polynomial system (Theorem~1 and Corollary~1).  For
  qubits, Eq.~(87) reads
  \begin{equation}
    M_{ba}=\frac{2b+1}{n+1}\,
    {}_4F_3\!\left(\begin{matrix}-b,\ b+1,\ -a,\ a+1\\ 1,\ -n,\ n+2\end{matrix};1\right),
    \qquad 0\leq a,b\leq n,
    \label{eq:2fba-macwilliams}
  \end{equation}
  where the generalized hypergeometric series terminates and $a,b$ label the
  spin sectors of the operator space.  Equation~\eqref{eq:2fba-macwilliams} makes
  the linear-programming bounds of the preceding item explicit for every $n$,
  but the preprint does not evaluate them for $t\geq2$
  \sourcecite{ref:2fba-teixeira}{Tei26}.
\end{itemize}

\subsection*{References}

\begin{enumerate}[label={},leftmargin=4.8em,labelsep=0.5em]
  \footnotesize
  \item[\textup{[PR04]}]\label{ref:2fba-pollatsek-ruskai}
  H. Pollatsek and M. B. Ruskai, ``Permutationally invariant codes for quantum
  error correction,'' \emph{Linear Algebra and its Applications} \textbf{392},
  255--288 (2004).
  \href{https://doi.org/10.1016/j.laa.2004.06.014}{doi:10.1016/j.laa.2004.06.014};
  \href{https://arxiv.org/abs/quant-ph/0304153}{arXiv:quant-ph/0304153}.

  \item[\textup{[Ouy14]}]\label{ref:2fba-ouyang}
  Y. Ouyang, ``Permutation-invariant quantum codes,'' \emph{Physical Review A}
  \textbf{90}, 062317 (2014).
  \href{https://doi.org/10.1103/PhysRevA.90.062317}{doi:10.1103/PhysRevA.90.062317};
  \href{https://arxiv.org/abs/1302.3247}{arXiv:1302.3247}.

  \item[\textup{[AAB24]}]\label{ref:2fba-aab}
  A. Aydin, M. A. Alekseyev, and A. Barg, ``A family of permutationally
  invariant quantum codes,'' \emph{Quantum} \textbf{8}, 1321 (2024).
  \href{https://doi.org/10.22331/q-2024-04-30-1321}{doi:10.22331/q-2024-04-30-1321};
  \href{https://arxiv.org/abs/2310.05358}{arXiv:2310.05358}.

  \item[\textup{[BMO+26]}]\label{ref:2fba-bond-et-al}
  L. J. Bond, J. Min\'a\v{r}, M. Ozols, A. Safavi-Naini, and V. Visnevskyi,
  ``Permutation-invariant codes: a numerical study and qudit constructions,''
  arXiv preprint (2026).
  \href{https://arxiv.org/abs/2603.10981}{arXiv:2603.10981}.

  \item[\textup{[AAB26]}]\label{ref:2fba-aab-convex}
  A. Aydin, V. V. Albert, and A. Barg, ``Quantum error correction beyond
  $\mathrm{SU}(2)$: spin, bosonic, and permutation-invariant codes from convex
  geometry,'' \emph{PRX Quantum} \textbf{7}, 010341 (2026).
  \href{https://doi.org/10.1103/kx3b-4nrp}{doi:10.1103/kx3b-4nrp};
  \href{https://arxiv.org/abs/2509.20545}{arXiv:2509.20545}.

  \item[\textup{[KT26]}]\label{ref:2fba-kubischta-teixeira}
  E. Kubischta and I. Teixeira, ``MacWilliams identities for intrinsic quantum
  codes,'' arXiv preprint, version 2 (2026).
  \href{https://arxiv.org/abs/2604.16023}{arXiv:2604.16023}.

  \item[\textup{[Tei26]}]\label{ref:2fba-teixeira}
  I. Teixeira, ``Orthogonal polynomials and the MacWilliams transform for
  permutation-invariant qudit codes,'' arXiv preprint (2026).
  \href{https://arxiv.org/abs/2605.15372}{arXiv:2605.15372}.
\end{enumerate}

\subsection*{Comment}

The case $t=1$ is settled: the seven-qubit codes of Pollatsek and Ruskai attain
Eq.~\eqref{eq:2fba-conjecture}, and the unrestricted exclusion of shorter codes
comes from the linear program of Kubischta and Teixeira, which is a preprint.
For $t=2$ the nineteen-qubit code of Aydin, Alekseyev, and Barg settles
existence at the conjectured length, so the remaining question is whether a
permutation-invariant distance-five code exists on at most eighteen qubits.  For
$t\geq3$ neither an exact code of length $3t^2+3t+1$ nor an unrestricted lower
bound of that order was located; the explicit constructions give only
Eq.~\eqref{eq:2fba-aab-bound}, and the numerical searches are not proofs.  A
proof of Eq.~\eqref{eq:2fba-conjecture} needs both a matching unrestricted lower
bound and existence at the conjectured length for every $t$, while a single $t$
with a different value refutes it.  The record
\href{https://qiqc-op.com/problem/op_a72722e455d85d55/}{Minimal permutation-invariant qubit code for two errors}
isolates the case $t=2$; resolving it would not decide
Eq.~\eqref{eq:2fba-conjecture} for other $t$.  Literature checked through
15 September 2026.

\subsection*{Field}

Quantum Error Correction

\subsection*{Topic}

Quantum coding theory

\subsection*{ID}

\texttt{op\_2fba72464e1b4de4}
